% 2007 年四川卷理科第 19 题：直角梯形 PCBM 与平面外的点 A。
% 来源 PDF：gaokaomath-src/普通高考/2007/2007四川理.pdf，SHA256 6f3287a6ec1c136fb7e457cb0264d069c44eb8ff58a81cc3598f9c8f1395cc2d
% 原页：2007四川理.pdf 第 2 页（A3 三栏）第 1 栏，书页码 410，图在第 19 题题干下方。
% 重绘依据：P、M 在上方，PM 水平（长 1）；C 在 P 正下方（PC=1，竖直），B 在 C 右侧
%   水平线上（BC=2），PM∥BC；M、B 连线向右下；A 在左下方（按原图取 A=P+(-0.95,-1.45)）。
%   可见线段 PM、PA、MB、AB、AM 画实线；被梯形遮住的 PC、CB、AC 画虚线；
%   C 处 PC 与 CB 之间画直角记号（∠PCB=90°）。
\begin{tikzpicture}[scale=2.2, semithick,
  every node/.style={inner sep=1.5pt}]
  \coordinate (P) at (0,3);
  \coordinate (M) at (1,3);
  \coordinate (C) at (0,2);
  \coordinate (B) at (2,2);
  \coordinate (A) at (-0.95,1.55);
  % 被遮挡的三条线
  \draw[dashed] (P) -- (C);
  \draw[dashed] (C) -- (B);
  \draw[dashed] (A) -- (C);
  % 可见线段
  \draw (P) -- (M);
  \draw (P) -- (A);
  \draw (M) -- (B);
  \draw (A) -- (B);
  \draw (A) -- (M);
  % C 处直角记号
  \draw ($(C)+(0.18,0)$) -- ($(C)+(0.18,0.18)$) -- ($(C)+(0,0.18)$);
  \node[anchor=south east] at ([xshift=-1pt,yshift=1pt]P) {$P$};
  \node[anchor=south west] at ([xshift=1pt,yshift=1pt]M) {$M$};
  \node[anchor=north west] at ([xshift=2pt,yshift=-2pt]C) {$C$};
  \node[anchor=north west] at ([xshift=1pt,yshift=-1pt]B) {$B$};
  \node[anchor=north east] at ([xshift=-1pt,yshift=-1pt]A) {$A$};
\end{tikzpicture}
